%!TEX root=../../note

\subsection{Mathematical Induction}
\begin{definition}[Principle of Mathematical Induction]
    Let $P(n)$ be a statement depending on $n\in\N$. If the following two statements are true:
    \begin{enumerate}[label=(\arabic*)]
        \item $P(1)$ is true;
        \item for every $k\in\N$, if $P(k)$ is true, then $P(k+1)$ is true;
    \end{enumerate}
    then
    \begin{equation*}
        \forall n\in\N,\quad P(n)
    \end{equation*}
    is true.
\end{definition}

\begin{remark}[How to use induction in practice]
    To prove that $P(n)$ holds for every $n\in\N$ by induction, it suffices to carry out two steps: first, verify the \emph{base case} $P(1)$ directly; second, carry out the \emph{inductive step} by fixing an arbitrary $k\in\N$, assuming $P(k)$ holds (the \emph{inductive hypothesis}), and using this assumption to deduce $P(k+1)$. Once both steps are complete, the principle of mathematical induction guarantees $P(n)$ for all $n\in\N$ — no further verification of individual cases is needed.
\end{remark}

Here is an example: 
\begin{proposition}[Sum of the First $n$ Natural Numbers]
    For every $n\in\N$,
    \begin{equation*}
        1+2+\cdots+n=\frac{n(n+1)}{2}.
    \end{equation*}
\end{proposition}


\begin{proof}
    Define $P(n)$ to be the statement
    \begin{equation*}
        1+2+\cdots+n=\frac{n(n+1)}{2}.
    \end{equation*}
    The base case is $n=1$:
    \begin{equation*}
        1=\frac{1(1+1)}{2}.
    \end{equation*}

    For the inductive step, let $k\in\N$ and assume that $P(k)$ is true. Then
    \begin{align*}
        1+2+\cdots+(k+1)
            &= \frac{k(k+1)}{2}+(k+1) \\
            &= \frac{k(k+1)+2(k+1)}{2} \\
            &= \frac{(k+1)(k+2)}{2}.
    \end{align*}
    Thus, $P(k+1)$ is true. By the principle of mathematical induction, the proposition follows.
\end{proof}

The principle of mathematical induction has the following equivalent set formulation.
\begin{proposition}
    Let $S\subseteq\N$. If
    \begin{enumerate}
        \item $1\in S$;
        \item for every $k\in S$, $k+1\in S$;
    \end{enumerate}
    then $S=\N$.
\end{proposition}

\begin{proof}
    Define $P(n)$ to be the statement $n\in S$. The two assumptions give the base case $P(1)$ and the inductive implication $P(k)\Rightarrow P(k+1)$. Hence, by mathematical induction, every $n\in\N$ belongs to $S$. Since $S\subseteq\N$, it follows that $S=\N$.
\end{proof}

\subsubsection{Well-Ordering Principle}
\begin{theorem}[Well-Ordering Principle]
    Every nonempty subset $S$ of $\N$ has a smallest element.
\end{theorem}

\begin{remark}[Proof idea]
    We argue by contradiction: suppose $S$ has no smallest element. We want to show this forces $S=\emptyset$, contradicting the hypothesis that $S$ is nonempty. The natural tool is induction, so we build a set $T$ designed to capture, for each $n$, whether $S$ has been "avoided" up through $n$; showing $T=\N$ by induction will then show $S$ is empty.
\end{remark}

\begin{proof}
    Let $S\subseteq\N$ be nonempty. Suppose, for contradiction, that $S$ has no smallest element. Define
    \begin{equation*}
        T=\set{n\in\N : \set{1,2,\ldots,n}\cap S=\emptyset}.
    \end{equation*}
    That is, $T$ consists of exactly those $n\in\N$ for which none of $1,2,\ldots,n$ lies in $S$ — informally, $n\in T$ means $S$ has "not yet been seen" by the time we reach $n$. If we can show $T=\N$, then in particular no natural number lies in $S$ (since every $n\in\N$ would satisfy $n\in T$, forcing $n\notin S$), giving $S=\emptyset$ and the desired contradiction. We show $T=\N$ by induction on $n$.

    \textbf{Base case.} Since $S$ has no smallest element and $1$ is the smallest element of $\N$, we must have $1\notin S$ (otherwise $1$ would be the smallest element of $S$). Hence $\set{1}\cap S=\emptyset$, so $1\in T$.

    \textbf{Inductive step.} Suppose $k\in T$, i.e., $\set{1,2,\ldots,k}\cap S=\emptyset$. We show $k+1\in T$. Suppose, for contradiction, that $k+1\notin T$; then $\set{1,2,\ldots,k+1}\cap S\neq\emptyset$. Since $\set{1,\ldots,k}\cap S=\emptyset$ by the inductive hypothesis, the only way for $\set{1,\ldots,k+1}\cap S$ to be nonempty is if $k+1\in S$. But then $k+1\in S$ while none of $1,2,\ldots,k$ lie in $S$, so $k+1$ would be the smallest element of $S$ — contradicting our assumption that $S$ has no smallest element. Hence $k+1\in T$.

    By mathematical induction, $T=\N$. This implies that no natural number belongs to $S$, so $S=\emptyset$, contradicting the assumption that $S$ is nonempty. Therefore, $S$ has a smallest element.
\end{proof}

\begin{remark}[Practical use of well-ordering]
    In practice, the well-ordering principle is most often invoked as follows: to show a property $P(n)$ holds for every $n\in\N$, argue by contradiction. If $P(n)$ failed for some $n$, the set of counterexamples $\set{n\in\N : \neg P(n)}$ would be a nonempty subset of $\N$, and the well-ordering principle would guarantee it has a least element $k$. One then derives a contradiction using the minimality of $k$ — typically by showing that $P(k-1)$ must hold (since $k-1$ is smaller than $k$ and hence not a counterexample) and that $P(k-1)$ forces $P(k)$, contradicting $k$ being a counterexample. This is exactly the strategy used in the proposition below.
\end{remark}

\begin{proposition}
    Using the well-ordering principle, prove that for every $n\in\N$,
    \begin{equation*}
        1^2 + 2^2 + \cdots + n^2 = \frac{n\paren{n + 1}\paren{2n + 1}}{6}.
    \end{equation*}
\end{proposition}

\begin{remark}[Proof idea]
    We prove the claim by minimal counterexample. Suppose the formula failed for some $n$; collect all such failures into a set $S=\set{n\in\N:\neg P(n)}$ — this is the natural set to consider because if we can rule out $S$ being nonempty, then $P(n)$ holds for every $n$. Since $S\subseteq\N$, the well-ordering principle (rather than induction directly) is the right tool to extract a concrete least counterexample $k=\min S$. We then show $P(k-1)\Rightarrow P(k)$ by direct computation, so $k$ cannot actually be a counterexample — a contradiction.
\end{remark}

\begin{remark}[Why $\min$, not $\sup$]
    It is worth pausing on why we take the \emph{minimum} of $S$ here rather than some supremum-type argument. The well-ordering principle applies specifically to nonempty subsets of $\N$ (not to general subsets of $\R$), and it guarantees the existence of a smallest element that actually \emph{belongs to} $S$ — this is what allows us to write $k=\min S\in S$ and conclude $\neg P(k)$. This is a stronger statement than merely having an infimum: for a general subset of $\R$, an infimum need not belong to the set, so we could not conclude $\neg P(\inf S)$. Here, because $S\subseteq\N$ and $S$ is nonempty, well-ordering hands us an actual least element $k$ with $\neg P(k)$ true, which is exactly what is needed to test $P$ at $k$ and at $k-1$.

    It is also worth noting that induction and the well-ordering principle are logically equivalent ways of expressing the same underlying fact about $\N$: the proof of the well-ordering principle above used induction, and here we use well-ordering to reprove a statement ($P(n)$ for all $n$) that could equally well have been proved by direct induction on $n$. The two techniques are simply different packagings of the same idea.
\end{remark}

\mdfsetup{nobreak=true}
\begin{proof}
    Let $P(n)$ denote the statement
    \begin{equation*}
        1^2 + 2^2 + \cdots + n^2 = \frac{n(n + 1)(2n + 1)}{6}.
    \end{equation*}
    Suppose, for contradiction, that $P(n)$ fails for some $n$; that is,
    \begin{equation*}
        \neg\paren{\forall n \in \N,\ P(n)}
        \quad\text{equivalently}\quad
        \exists n \in \N,\ \neg P(n).
    \end{equation*}
    Then the set
    \begin{equation*}
        S := \set{n \in \N : \neg P(n)}
    \end{equation*}
    is nonempty. By the well-ordering principle, $S$ has a smallest element $k$. Since $P(1)$ holds, $1\notin S$, so $k > 1$ and hence $k - 1\in\N$.

    By minimality of $k$, we have $k - 1\notin S$, so $P(k - 1)$ is true:
    \begin{equation*}
        1^2 + 2^2 + \cdots + (k - 1)^2
            = \frac{(k - 1)k\paren{2(k - 1) + 1}}{6}
            = \frac{(k - 1)k(2k - 1)}{6}.
    \end{equation*}
    Adding $k^2$ to both sides gives
    \begin{align*}
        1^2 + \cdots + (k - 1)^2 + k^2
            &= \frac{(k - 1)k(2k - 1)}{6} + k^2 \\
            &= \frac{(k - 1)k(2k - 1) + 6k^2}{6} \\
            &= \frac{k\bigl((k-1)(2k-1)+6k\bigr)}{6} \\
            &= \frac{k(2k^2+3k+1)}{6}
             = \frac{k(k+1)(2k+1)}{6},
    \end{align*}
    so $P(k)$ is in fact true. This contradicts $k\in S$. Therefore $S=\emptyset$, and $P(n)$ holds for every $n\in\N$.
\end{proof}
\mdfsetup{nobreak=false}

\begin{remark}[General pattern: minimal counterexample via well-ordering]
    To show $P(n)$ holds for all $n\in\N$: assume the counterexample set $S=\set{n:\neg P(n)}$ is nonempty, let $k=\min S$ (which exists by well-ordering), show $P(k-1)$ holds by minimality of $k$, then show $P(k-1)\Rightarrow P(k)$ — contradicting $k\in S$.
\end{remark}
